\addtocontents{toc}{\protect\Needspace{3\baselineskip}}
\originalchapter{经典弦与几何习题}
本章十题为 8.251 自足公开作业题, 解答由 AI 独立推导. 其他完整公开重述题见后续习题章; 只有教材题号的项目仍列缺失, 不按编号补造题面.

\src{8.251 Spring 2007, HW1.2, (a)--(c), 10 分.}
\begin{exercise}
Lorentz 矩阵满足 $L^T\eta L=\eta$, $\eta=\operatorname{diag}(-1,1,\ldots,1)$.
(a) 证明两 Lorentz 矩阵乘积仍满足条件; (b) 证明逆矩阵满足条件; (c) 证明转置矩阵满足条件.
\end{exercise}
\begin{solution}
(a) $(L_1L_2)^T\eta(L_1L_2)=L_2^T(L_1^T\eta L_1)L_2=\eta$.
(b) 行列式给 $(\det L)^2=1$, 故可逆. 左右乘逆矩阵得 $(L^{-1})^T\eta L^{-1}=\eta$.
(c) 原式也给 $L^{-1}=\eta L^T\eta$. 对原式取逆并用 $\eta^{-1}=\eta$ 得 $L\eta L^T=\eta$, 正是 $L^T$ 的 Lorentz 条件. 这涉及全 Lorentz 群; 若限定正定向、保时向的分支, 还须单独核查相应符号条件.
\end{solution}

\src{8.251 HW1.5, 10 分.}
\begin{exercise}
$M,N\ge2$, $M<N$, 互素. 对 $z=x+\ii y$ 施加 $z\sim\ee^{2\pi\ii M/N}z$. 说明 $0\le\arg z<2\pi M/N$ 为什么通常不是基本域, 找真正基本域. 可用 Bézout 恒等式; 原题给练习 $187m+35n=1$.
\end{exercise}
\begin{solution}
互素使存在整数 $m,n$ 满足 $mM+nN=1$, 所以生成元的 $m$ 次幂就是 $\ee^{2\pi\ii/N}$. 生成的轨道与标准 $\mathbb Z_N$ 旋转群相同, 每个非原点轨道恰有 $N$ 个代表. 因而基本域可取 $0\le\arg z<2\pi/N$, 两边识别, 原点为固定点. 较大扇形包含 $M$ 份代表; 例如 $M=2,N=3$ 的 $4\pi/3$ 扇形会重复收录同一轨道. 欧几里得算法给 $1=3\cdot187-16\cdot35$, 可验证提示中的算术例子.
\end{solution}

\src{8.251 HW2.3, 10 分.}
\begin{exercise}
参考系 $S$ 中边长 $L$ 的静止立方体均匀带电, 固有密度 $\rho_0$, 电流为零. $S'$ 沿 $+x$ 以速度 $v$ 相对 $S$ boost.
用电荷不变性求 $\rho',\boldsymbol j'$, 验证 $(c\rho,\boldsymbol j)$ 是四矢量, 并用固有密度和四速度写电流公式, 检查两系.
\end{exercise}
\begin{solution}
$S'$ 同时面上的体积是 $L^3/\gamma$, 总电荷 $Q=\rho_0L^3$ 不变, 故 $\rho'=\gamma\rho_0$. 盒子在 $S'$ 中速度为 $-v\hat x$, $j'^x=-\gamma\rho_0v$, 其余分量为零. Boost 将 $(c\rho_0,0)$ 变为 $(\gamma c\rho_0,-\gamma v\rho_0)$, 与结果相符. 若四速度定义为 $u^\mu=\dd x^\mu/\dd\tau$, 则 $j^\mu=\rho_0u^\mu$. $S$ 中 $u=(c,0)$, $S'$ 中 $u'=(\gamma c,-\gamma v,0,0)$. 密度变化来自同时面与长度收缩, 不是电荷本身改变.
\end{solution}

\src{8.251 HW2.4, 10 分.}
\begin{exercise}
$S$ 中静止的大平行板电容器两板面电荷 $\pm\sigma_0$, $S'$ 沿 $+x$ boost.
(i) 板为定 $y$ 平面, 场 $E_\perp$ 沿 $y$; (ii) 板为定 $x$ 平面, 场 $E_\parallel$ 沿 $x$.
由电荷不变性求两场变换规律, 再用 $F'^{\mu\nu}=L^\mu{}_{\rho}L^\nu{}_{\sigma}F^{\rho\sigma}$ 核查. $S$ 中只含 $E_x,E_y$.
\end{exercise}
\begin{solution}
(i) 板面沿 $x$ 的长度收缩, 面积减为 $1/\gamma$, 故 $\sigma'=\gamma\sigma_0$, Gauss 关系给 $E'_y=\gamma E_y$.
(ii) 板面位于 $yz$ 方向, 面积不变, $\sigma'=\sigma_0$, 所以 $E'_x=E_x$. 两种情况板间距离的改变不控制无限板场强, 面电荷才控制它.

按原课有理化单位选 $F^{0i}=E_i$, $F^{ij}=\epsilon^{ijk}B_k$, 非零 boost 分量 $L^0{}_0=L^x{}_x=\gamma$, $L^0{}_x=L^x{}_0=-\gamma\beta$. 对 $F^{0x}$ 的二项组合给 $(\gamma^2-\gamma^2\beta^2)E_x=E_x$, 对 $F^{0y}$ 给 $\gamma E_y$. 同时 $F'^{xy}=-\gamma\beta E_y$, 所以还产生 $B'_z=-\gamma vE_y/c$. 只变换电场而声称磁场仍零, 不能满足同一个张量变换.
\end{solution}

\src{8.251 HW3.4, (a)--(c), 10 分.}
\begin{exercise}
非相对论闭弦绕无摩擦圆柱绷紧, 圆周 $x\sim x+2\pi R$, 轴向坐标 $y$. 张力 $T_0$, 线密度 $\mu_0$, $v_0=\sqrt{T_0/\mu_0}$, 通解 $y=h_+(x-v_0t)+h_-(x+v_0t)$.
(a) 写周期条件并证明 $h_+',h_-'$ 周期.
(b) 证明可写 $h_+(u)=\alpha u+f(u),h_-(v)=\beta v+g(v)$, $f,g$ 周期, 求 $\alpha,\beta$ 关系.
(c) 求轴向总动量, 是否守恒?
\end{exercise}
\begin{solution}
记 $L=2\pi R$. 对任意独立 $u=x-v_0t,v=x+v_0t$, 周期性给 $h_+(u+L)-h_+(u)=-[h_-(v+L)-h_-(v)]$. 左右只依各自独立变量, 必为常数. 对 $u,v$ 求导证明两导数周期. 令常数为 $\alpha L$ 与 $\beta L$, 则 $\alpha+\beta=0$, 减去线性项后的 $f,g$ 都以 $L$ 为周期.

轴向速度 $\dot y=-v_0h_+'+v_0h_-'$. 乘 $\mu_0$ 沿一周积分, 周期函数的导数积分为零, 得 $P_y=\mu_0Lv_0(\beta-\alpha)=-2\mu_0Lv_0\alpha$, 与时间无关. 线性项不表示闭弦沿周方向不单值; 它们的 $x$ 部分相消, 剩下整体轴向匀速平移.
\end{solution}

\src{8.251 HW4.3, (a)--(b), 10 分.}
\begin{exercise}
二维实向量空间有常度规, 存在类空 $s'$ 和类时 $t'$.
(a) 构造单位 $s,t$ 满足 $s^2=1,t^2=-1,s\cdot t=0$, 给 $\mathbb R^{1,1}$ 标准基.
(b) 对任意 $v_1,v_2$ 证明 $(v_1\cdot v_2)^2\ge v_1^2v_2^2$, 说明取等.
\end{exercise}
\begin{solution}
先取 $t=t'/\sqrt{-t'^2}$, 再令 $s_0=s'+(s'\cdot t)t$. 它与 $t$ 正交, 且 $s_0^2=s'^2+(s'\cdot t)^2>0$, 所以 $s=s_0/\sqrt{s_0^2}$ 合条件. 任意 $v_i=a_it+b_is$, 内积为 $-a_1a_2+b_1b_2$, 展开后
\[
(v_1\cdot v_2)^2-v_1^2v_2^2=(a_1b_2-a_2b_1)^2\ge0.
\]
等号当两向量线性相关, 包括某个向量为零. 这项反向 Schwarz 型不等式依二维 Lorentz 度规, 不可以不加条件直接搬到一般高维任意向量.
\end{solution}

\src{8.251 HW4.4, (a)--(d), 10 分.}
\begin{exercise}
单位球 $x^2+y^2+z^2=1$ 从北极 $N=(0,0,1)$ 投影到 $z=0$ 平面参数 $(\xi_1,\xi_2)$.
(a) 画构造, 给覆盖除北极外整球所需参数范围.
(b) 求 $x,y,z$; (c) 求诱导度规四分量; (d) 积分检验球面积.
\end{exercise}
\begin{solution}
参数范围是 $\mathbb R^2$. 连接线写为 $(t\xi_1,t\xi_2,1-t)$, 与球面的非北极交点满足 $t=2/(1+q)$, $q=\xi_1^2+\xi_2^2$. 因而 $x=2\xi_1/(1+q),y=2\xi_2/(1+q),z=(q-1)/(q+1)$. 平面无穷远逼近北极.
\begin{center}\begin{tikzpicture}[scale=1.7]
\draw(0,0)circle(1);\draw[-{Stealth}](-1.3,0)--(2.3,0)node[right]{$z=0$};
\coordinate(N)at(0,1);\coordinate(P)at(2,0);\coordinate(Q)at(.8,.6);
\draw[structurecolor](N)--(P);\fill(N)circle(.025)node[above]{$N$};\fill(P)circle(.025)node[below]{$(\xi,0)$};\fill(Q)circle(.025)node[right]{$(x,z)$};
\end{tikzpicture}\end{center}
图是 $y=0$ 截面, 示意取 $\xi=2$; 不把截面图当全三维投影. 对坐标导数逐项内积, 混合项相消, 得 $g_{ij}=4\delta_{ij}/(1+q)^2$, 即 $g_{11}=g_{22}=4/(1+q)^2,g_{12}=g_{21}=0$. 面积元 $\sqrt g=4/(1+q)^2$, 极坐标积分 $A=8\pi\int_0^\infty r\dd r/(1+r^2)^2=4\pi$. 单个未覆盖北极面积为零, 不影响积分.
\end{solution}

\src{8.251 HW5.2, 5 分.}
\begin{exercise}
四维时空一静止 D1 膜位于 $x^3=0$, 相对 $x^1$ 轴逆时针旋转 $\theta$. 给开弦两端的完整边界条件, 用 $X^\mu(\tau,\sigma_*)$ 与世界面边界动量流 $P^\sigma_\mu$ 表示.
\end{exercise}
\begin{solution}
膜切向 $e_\parallel=(\cos\theta,\sin\theta,0)$, 平面内法向 $e_\perp=(-\sin\theta,\cos\theta,0)$. 每个端点分别需 $P^\sigma_0=0$, $\cos\theta P^\sigma_1+\sin\theta P^\sigma_2=0$, $-\sin\theta X^1+\cos\theta X^2=0$, $X^3=0$. 后两条是固定在穿原点膜上的 Dirichlet 条件, 同时其 $\tau$ 导数为零. 如果膜平移, 法向固定值应换成其位置; 本题未给平移.
\end{solution}

\src{8.251 HW6.1, (a)--(d), 15 分, 官方修正版 2007年3月17日.}
\begin{exercise}
闭弦 $\boldsymbol X=\frac12[\boldsymbol F(u)+\boldsymbol G(v)]$, $u=ct+\sigma,v=ct-\sigma$, $\sigma\sim\sigma+L$, $L=E/T_0$. 两函数以 $L$ 为周期, $|F'|=|G'|=1$. 在 $u_0,v_0$ 处 $F'=G'$, 单位球上的两条导数路径正则横交, 即 $F'',G''$ 非零且不平行.
(a) 平移坐标使 cusp 为原点, 推导 $X=T(\sigma-\sigma_0)^2+R(\sigma-\sigma_0)^3+\cdots$, $T=(F''+G'')/4,R=(F'''-G''')/12$, 说明 $T\ne0,T\perp F'$, 特殊时 $R$ 可零.
(b) 以正 $y$ 轴沿 $T$, $x$ 轴使 $R$ 在 $xy$ 平面, 说明普通 cusp 的 $y\sim x^{2/3}$, 速度位于什么平面?
(c) 取 $F=(L/2\pi)(\sin(2\pi u/L),-\cos(2\pi u/L),0)$, $G=(L/4\pi)(\sin(4\pi v/L),0,-\cos(4\pi v/L))$, 验条件, 给 $t=\sigma=0$ cusp 的方向、平面、速度并画示意.
(d) 证明整体形状周期 $L/(4c)$, 每周期有几个 cusp?
\end{exercise}
\begin{solution}
固定 $t=t_0$ 令 $s=\sigma-\sigma_0$, 两自变量增量为 $s,-s$, 常数项平移去掉, 一次项因 $F'=G'$ 消失, 直接 Taylor 展开得 (a). $F'\cdot F''=G'\cdot G''=0$ 使 $T$ 垂直共同切向; 横交使 $F'',G''$ 非平行且非零, 它们之和不能为零. $R$ 的公式含更高导数, 不由横交单独保证非零.

若 $R$ 在 $T$ 正交方向的投影非零, 则局部 $y=|T|s^2+O(s^3)$、$x=R_xs^3+O(s^4)$, 消去 $s$ 得 $y\propto|x|^{2/3}$. 这是普通 cusp 的附加非退化条件; 若 $R=0$ 或平行 $T$, 必须查看更高阶, 原题的幂律不能无条件宣称. 速度 $\dot X=cF'$ 位于垂直 $T$ 的平面, 一般不必在 $\operatorname{span}(T,R)$ 内.

例中令 $k=2\pi/L$, $F'=(\cos ku,\sin ku,0)$、$G'=(\cos2kv,0,\sin2kv)$, 满足单位模和周期. 在零点共同方向是 $+x$, 位置为 $(0,-L/4\pi,-L/8\pi)$, 可平移去掉. $T=(0,k/4,k/2)$、$R=(k^2/4,0,0)$, cusp 形状平面为 $\operatorname{span}[(1,0,0),(0,1,2)]$, 速度 $(c,0,0)$. 示意只画该平面中的标准局部形式, 横纵轴按 $R$ 和 $T$ 方向, 并非全弦全局曲线:
\begin{center}\begin{tikzpicture}\begin{axis}[width=.6\linewidth,height=4.5cm,xlabel={$R$ 方向},ylabel={$T$ 方向},axis lines=middle,xtick=\empty,ytick=\empty]
\addplot[structurecolor,domain=-1:1,samples=101]({x^3},{x^2});
\end{axis}\end{tikzpicture}\end{center}
时间增 $L/(4c)$ 并重新标号 $\sigma\to\sigma-L/4$ 时, $u$ 不变、$v$ 增 $L/2$, $G$ 已重复, 故几何形状相同. 因此 $L/(4c)$ 是本例的形状周期. cusp 条件 $u=mL/2,v=nL/4$, $m,n$ 同奇偶. 在 $0\le t<L/(4c)$、$0\le\sigma<L$ 内, 分别有 $(t,\sigma)=(0,0)$ 与 $(L/(8c),3L/8)$ 两个事件, 故每形状周期两个 cusp. 固定参数标签的重复周期与形状重复周期不同.
\end{solution}

\src{8.251 HW6.2, 5 分.}
\begin{exercise}
宇宙弦产生圆锥亏角 $\Delta$. 观察者至弦距离 $d$, 类星体至弦距离 $\ell$, 其位置使出现双像. 在小亏角近似下求双像角分离 $\delta\phi$.
\end{exercise}
\begin{solution}
沿弦垂直截面展开圆锥, 相邻两份源的横向间隔约为 $\ell\Delta$. 观察者至源总距离为 $d+\ell$, 因而 $\delta\phi\simeq\Delta\ell/(d+\ell)$. 这取弦垂直视线且两距离是同一平直小角几何的距离; 在宇宙学传播中应改用透镜到源与观察者到源的角直径距离比, 若弦倾斜还需相应投影. 亏角几何的双像机制与点质量引力偏折不同.
\end{solution}
